\documentclass{article}
\usepackage{danetiska-instagram}
\usepackage[noglossaries]{leipzig}
\begin{document}

\newcommand{\myleip}[1]{\dahead{\leipzigfont{#1}}}

% SAMPLE ONLY: replace the illustrative endings and captions with verified
% Danetian paradigms before publishing.
\begin{dacard}{Personal pronouns}{Danetian grammar / 01}
	\dadesc{First and second person singular pronouns.}
	\begin{datable}{KXX}
		\dahead{Case} & \myleip{1.sg} & \myleip{2.sg} \\
		\midrule
		Nominative	& \daform{edja} & \daform{tu} \\
		Accusative	& \daform{me} & \daform{tve} \\
		Genitive	& \daform{mene} & \daform{teve} \\
		Ablative	& \daform{med} & \daform{tved} \\
		Dative	& \daform{medjo} & \daform{tebio} \\
		Locative	& \daform{moi} & \daform{toi} \\
		Instrumental	& \daform{moi}	& \daform{toi} \\
	\end{datable}
\end{dacard}

\begin{dacard}{Personal pronouns}{Danetian grammar / 02}
	\dadesc{First and second person plural pronouns.}
	\begin{datable}{KXX}
		\dahead{Case} & \myleip{1.pl} & \myleip{2.pl} \\
		\midrule
		Nominative	& \daform{vi} & \daform{ju} \\
		Accusative	& \daform{nos} & \daform{vos} \\
		Genitive	& \daform{nysom} & \daform{usom} \\
		Ablative	& \daform{nysmos} & \daform{usmos} \\
		Dative	& \daform{nysmos} & \daform{usmos} \\
		Locative	& \daform{nysu} & \daform{usu} \\
		Instrumental	& \daform{nysbi}	& \daform{usbi} \\
	\end{datable}
\end{dacard}



\begin{dacard}{Personal pronouns}{Danetian grammar / 03}
	\dadesc{Third person masculine pronouns, both singular and plural. Similar to the o-stem declension, but with major differences.}
	\begin{datable}{KXX}
		\dahead{Case} & \myleip{3.sg.m} & \myleip{3.pl.m} \\
		\midrule
		Nominative	& \daform{enos} & \daform{enoi} \\
		Accusative	& \daform{enom} & \daform{enons} \\
		Genitive	& \daform{enosio} & \daform{enoisom} \\
		Ablative	& \daform{enosmod} & \daform{enoimos} \\
		Dative	& \daform{enosmoi} & \daform{enoimos} \\
		Locative	& \daform{enosmi} & \daform{enoisu} \\
		Instrumental	& \daform{enosme}	& \daform{enois} \\
	\end{datable}
\end{dacard}


\begin{dacard}{Personal pronouns}{Danetian grammar / 04}
	\dadesc{Third person feminine pronouns, both singular and plural. Similar to the a-stem declension, but with major differences.}
	\begin{datable}{KXX}
		\dahead{Case} & \myleip{3.sg.f} & \myleip{3.pl.f} \\
		\midrule
		Nominative	& \daform{ena} & \daform{enai} \\
		Accusative	& \daform{enam} & \daform{enans} \\
		Genitive	& \daform{enasias} & \daform{enasom} \\
		Ablative	& \daform{enasmad} & \daform{enamos} \\
		Dative	& \daform{enasmai} & \daform{enamos} \\
		Locative	& \daform{enasmi} & \daform{enasu} \\
		Instrumental	& \daform{enasme}	& \daform{enabi} \\
	\end{datable}
\end{dacard}

\begin{dacard}{Personal pronouns}{Danetian grammar / 05}
	\dadesc{Third person neuter pronouns, both singular and plural. Similar to the o-stem declension, but with major differences.}
	\begin{datable}{KXX}
		\dahead{Case} & \myleip{3.sg.n} & \myleip{3.pl.n} \\
		\midrule
		Nominative	& \daform{enod} & \daform{ena} \\
		Accusative	& \daform{enod} & \daform{ena} \\
		Genitive	& \daform{enosio} & \daform{enoisom} \\
		Ablative	& \daform{enosmod} & \daform{enoimos} \\
		Dative	& \daform{enosmoi} & \daform{enoimos} \\
		Locative	& \daform{enosmi} & \daform{enoisu} \\
		Instrumental	& \daform{enosme}	& \daform{enois} \\
	\end{datable}
\end{dacard}



\end{document}
